\documentclass[10pt,a4paper]{amsart}
\usepackage[margin=19mm]{geometry}
\usepackage{amsmath,amssymb,mathtools}
\usepackage[scr=rsfs,cal=euler]{mathalfa}
\usepackage{xfrac}
\usepackage[unicode,hidelinks,bookmarks=false]{hyperref}
\newcommand{\ie}{i.e.\ }
\newcommand{\cf}{cf.\ }
\newcommand{\resp}{respectively\ }
\newcommand{\Subsection}{\subsection}
\providecommand{\nobreakdash}{\nobreak\hbox{-}}
\setlength{\emergencystretch}{2em}
\multlinegap0pt
\usepackage{amssymb}
\usepackage{xcolor}
\usepackage[shortlabels]{enumitem}
\newtheorem{definition}{Definition}
\newtheorem{lemma}{Lemma}
\newcommand{\ai}{A}
\newcommand{\bb}{B}
\newcommand{\cc}{C}
\title{Potential Carroll Structures and Special Carrollian Manifolds}
\author{Samuel Blitz, Gabriel Herczeg, David McNutt}
\date{Source: arXiv:2601.20068v1 (CC BY 4.0)}
\begin{document}
\maketitle

\noindent\emph{Source and licence.} Potential Carroll Structures and Special Carrollian Manifolds. Version arXiv:2601.20068v1 (CC BY 4.0), distributed by arXiv under the Creative Commons Attribution 4.0 International licence, http://creativecommons.org/licenses/by/4.0/, as recorded in the arXiv licence metadata for this exact version and shown on the abstract page https://arxiv.org/abs/2601.20068v1. Full licence notice: \texttt{licenses/arxiv-2601.20068-CC-BY-4.0.txt}.\par\smallskip

\noindent\textit{Editorial scope.} Counted unit: the article's lemma nonvanishing-torsion-lemma (source0.tex lines 229--240). The article's complete original proof of it is delivered with it (source0.tex lines 242--347, one \texttt{proof} environment of 106 lines). No mathematics is written, repaired or re-derived: the statement and the proof are the article's own lines, in the article's own notation, and no retained line is substituted or re-flowed.  Retained uncounted as the source-local context the proof needs: lines 180--188.  Nothing was omitted from any retained span, so the package carries no omission record.  No retained line contains a citation, so the wrapper carries no bibliography and no bibliography entry.  No retained line contains a figure environment, an image-inclusion command or drawing code, so no image is delivered and the package declares no asset dependency.  Editorial formatting only, in addition: an AMS \texttt{amsart} preamble that restates the article's own declarations and macros used by the retained lines, namely the environments \texttt{definition}, \texttt{lemma} on separate counters, together with the standard packages those lines need.  The retained environments print in this document as Definition 1, Lemma 1 in order of appearance, whereas the article prints the same items with the numbering of its own full document.  The complete native source of the article as distributed in the arXiv e-print for this exact version is bundled unchanged as \texttt{sources/193502/source0.tex}.
\begin{definition} \label{minimal-torsion-definition}
    Let $(M,g,\ell)$ be a Carrollian structure equipped with an Ehresmann connection $\omega$ such that $\omega(\ell) = 1$. Suppose $T \in \Gamma(TM \otimes \wedge^2 T^* M)$ is a section satisfying
        \begin{align}
        g(T(\ell,X), Y) &= \tfrac{1}{2} (\mathcal{L}_{\ell} g)(X,Y), ~~\forall X,Y \in \Gamma(T M)\,, \label{lie-g-constraint}\\
        \omega(T(\ell,X)) &= (\mathcal{L}_{\ell} \omega)(X) \,, \forall X,Y \in \Gamma(TM) \label{minimal-constraint} \\
        T(U,V) &= 0,~~\forall U,V \in \ker(\omega)\label{horiz-constraint}\,.
    \end{align}
    Then we say that the section $T$ is \textit{minimal}.
\end{definition}

\begin{lemma} \label{nonvanishing-torsion-lemma}
    Let $(M^3,g,\ell)$ be a Carrollian structure equipped with a linear connection $\nabla$ satisfying $\nabla \ell = 0 = \nabla g$.
    Suppose that there exists an Ehresmann connection $\omega$ such that the torsion of $\nabla$ is minimal with respect to the tuple $(M,g,\ell,\omega)$. Assuming the torsion is non-vanishing, there exists a symmetric tensor $N$ such that
    $$\nabla_{(a} \omega_{b)} = N_{ab}\qquad \text{ where } \qquad N(\ell, \cdot) = 0$$
    if and only if one of the following conditions holds:
        \begin{itemize}
        \item for $V := Trace(T)$,  $V(\ell) \neq 0$, 
        $$\nabla_{(a}\gamma_{b)} = N_{ab}, \qquad \gamma_b = \frac{V_b-\mathcal{L}_\ell \omega_b}{V(\ell)} \,, $$
        \item $V(\ell) = 0$  
        $$-\nabla_d T_{abc} = \nabla_d (\mathcal{L}_\ell g_{a[b}) \omega_{c]} + \mathcal{L}_\ell g_{a[b} N_{c]d}\,. $$
    \end{itemize}
\end{lemma}

\begin{proof}
    Choosing coordinates $(u,x^1,x^2)$ where $\ell = \partial_u$, without loss of generality, we may write any such Ehresmann connection as
    \begin{equation} 
        \omega = du + \omega_i (u,x,y) dx^i,\quad i =1,2. 
    \end{equation} 
    As the degenerate metric can be treated as a Riemannian signature metric on a 2-manifold, it is always possible to diagonalize this with an upper-triangular coframe:
     \begin{equation} \begin{aligned}
       m^1 &= m^1_{~1}(u, x^i) dx^1\,, \\
       m^2 &= m^2_{~1}(u,x^i) dx^1 + m^2_{~2}(u,x^i) dx^2.
    \end{aligned} \end{equation}
    Choosing $(\theta^\ai) = (\omega, m^1, m^2)$ as a coframe, the dual frame $(e_\ai)$ is then 
    \begin{equation} \begin{aligned}
        e_1 &= \ell = \partial_u, \\
        e_2 &= \frac{\omega_2 m^2_{~1} - \omega_1 m^2_{~2}}{m^1_{~1} m^2_{~2}} \partial_u + \frac{1}{m^1_{~1}} \partial_{x^1} - \frac{m^2_{~1}}{m^1_{~1} m^2_{~2}} \partial_{x^2}, \\
        e_3 &= - \frac{\omega_2}{m^2_{~2}} \partial_{u} + \frac{1}{m^2_{~2}} \partial_{x^1} .\\ 
    \end{aligned} \end{equation}
    In the frame basis, we can constrain the connection coefficients $\Gamma^\ai_{~\bb\cc}$ for $\nabla$ using the hypothesis. The conditions $\nabla \ell = \nabla g  =0$ yield purely algebraic expressions:
    \begin{equation}
    \begin{aligned}
        & \nabla \ell = 0 \rightarrow  \Gamma^{1}_{~1\ai}= \Gamma^{2}_{~1\ai}=\Gamma^{3}_{~1\ai} = 0, \quad \ai = 1,2,3, \\
        &\nabla g= 0 \rightarrow \Gamma^{2}_{~2\ai}=\Gamma^{3}_{~3\ai}=0, \Gamma^{3}_{~21} = -\Gamma^{2}_{~31}, \Gamma^{3}_{~22} = -\Gamma^{2}_{~32},\Gamma^{3}_{~23} = -\Gamma^{2}_{~33} 
    \end{aligned}
    \end{equation}
    giving
    \begin{equation} \begin{aligned}
        \nabla e_1 =& 0, \\
        \nabla e_2 =& \Gamma^1_{~21} e_1 \otimes \theta^1 + \Gamma^1_{~22} e_1 \otimes \theta^2 + \Gamma^1_{~23} e_1 \otimes \theta^3 \\ 
        & -\Gamma^2_{31} e_3 \otimes \theta^1 - \Gamma^2_{~32} e_3 \otimes \theta^2 - \Gamma^2_{~33} e_3 \otimes \theta^3 \\
        \nabla e_3 =& \Gamma^1_{~31} e_1 \otimes \theta^1 + \Gamma^1_{~32} e_1 \otimes \theta^2 + \Gamma^1_{~33} e_1 \otimes \theta^3  \\ 
        &+ \Gamma^2_{31} e_2 \otimes \theta^1 + \Gamma^2_{~32} e_2 \otimes \theta^2 + \Gamma^2_{~33} e_2 \otimes \theta^3.
    \end{aligned} \label{eqn:Thm32a} \end{equation}
    \noindent The covariant derivative of the Ehresmann connection is 
    \begin{equation} \begin{aligned}
        \nabla \omega = \nabla \theta^1 =& -\Gamma^1_{21} \theta^2 \otimes \theta^1 - \Gamma^1_{22} \theta^2 \otimes \theta^2 - \Gamma^1_{23} \theta^2 \otimes \theta^3 \\ 
        &- \Gamma^1_{31} \theta^3 \otimes \theta^1 - \Gamma^1_{32} \theta^3 \otimes \theta^2 - \Gamma^1_{33} \theta^3 \otimes \theta^3.  
    \end{aligned} \label{eqn:Thm32b} \end{equation}
    
    By hypothesis, $\omega$ is such that the torsion of $\nabla$ is minimal with respect to $(M,g,\ell,\omega)$. According to Definition~\ref{minimal-torsion-definition}, we may further determine $\Gamma$, using
    \begin{equation} \label{lie-derivs}
    \begin{aligned}
        \mathcal{L}_{\ell}g &= (\ln m^1_{~1})_{,u} \theta^2 \theta^2 + (\ln m^2_{~2})_{,u} \theta^3 \theta^3 + \\
        &\quad \quad 2 \left(\frac{(m^2_{~1})_{,u} m^2_{~2} - (m^2_{~2})_{,u} m^2_{~1}}{m^1_{~1} m^2_{~2}} \right) \theta^2 \theta^3, \\
        \mathcal{L}_\ell \omega &= \frac{\omega_{1,u} m^2_{~2} - \omega_{2,u} m^2_{~1}}{m^1_{~1} m^2_{~2}} \theta^2 + \frac{\omega_{2,u}}{m^2_{~2}} \theta^3.
    \end{aligned}
    \end{equation}
    Computing $T^\ai_{~\bb \cc}\ell^\cc = T(\ell, \cdot)$ using Cartan's first structure equation, we find
    \begin{equation}
    \begin{aligned}
        T(\ell,\cdot) &= \left( \Gamma^1_ {~21} + \frac{ \omega_{1,u} m^2_{~2}-\omega_{2,u}m^2_{~1}}{m^1_{~1} m^2_{~2}} \right) e_1 \otimes \theta^2 + \left(  \Gamma^1_ {~31} +\frac{\omega_{2,u}}{m^2_{~2}} \right) e_1 \otimes \theta^3 \\
        &\quad \quad  (\ln m^1_{~1})_{,u} e_2 \otimes \theta^2 + (\ln m^2_{~2})_{,u} e_3 \otimes \theta^3 + \Gamma^2_{~31} e_2 \otimes \theta^3 \\
        &\quad \quad -\left( \Gamma^2_{~31} +\frac{(m^2_{~1})_{,u} m^2_{~2} - (m^2_{~2})_{,u} m^2_{~1}}{m^1_{~1} m^2_{~2}} \right) e_2 \otimes \theta^3.
    \end{aligned}
    \end{equation}
    Combining the above with Definition~\ref{minimal-torsion-definition},
    we have the following expressions for connection coefficients: 
    \begin{equation}
    \begin{aligned}
        \Gamma^1_{~21} =& 0 \\ 
        \Gamma^1_{~31} =& 0, \\
         \Gamma^2_ {~31} =& -\left(\frac{m^2_{~2,u} m^2_{~1}- m^2_{~1,u} m^2_{~2} }{m^1_{~1} m^2_{~2}}\right), \\
         \Gamma^1_{~23} =& \Gamma^1_{~32} + \frac{ -\omega_{2,u} \omega_1 + \omega_{1,u} \omega_2  + \omega_{2,1}-\omega_{1,2}}{m^1_{~1} m^2_{~2}}, \\
         \Gamma^2_{~32} = & \frac{(m^1_{~1})_{,u} \omega_2 - (m^1_{~1})_{,2}}{m^1_{~1} m^2_{~2}}, \\
       \Gamma^2_{~33} = & - \left(\frac{(m^2_{~1})_{,u} \omega_2 - (m^2_{~2})_{,u} \omega_1 + (m^2_{~2})_{,1}-(m^2_{~1})_{,2}}{m^1_{~1} m^2_{~2}} \right).
    \end{aligned} \label{eqn:Thm32c}
    \end{equation}
    
    In light of this, to prove the first two conditions in the theorem, we express the torsion tensor explicitly: 
        \begin{equation}
    \begin{aligned}
        {\bf T} &= 2 \left( \frac{ \omega_{1,u} m^2_{~2}-\omega_{2,u}m^2_{~1}}{m^1_{~1} m^2_{~2}} \right) e_1 \otimes \theta^1 \wedge \theta^2 + 2\left( \frac{\omega_{2,u}}{m^2_{~2}} \right) e_1 \otimes \theta^1 \wedge \theta^3 \\ 
        &\quad + 2(\ln m^1_{~1})_{,u} e_2 \otimes \theta^1 \wedge \theta^2 + 2(\ln m^2_{~2})_{,u} e_3\otimes \theta^1  \wedge \theta^3  \\
        &\quad  + \quad 4 \left(\frac{(m^2_{~1})_{,u} m^2_{~2} - (m^2_{~2})_{,u} m^2_{~1}}{m^1_{~1} m^2_{~2}} \right) (e_3 \otimes \theta^1 \wedge \theta^2 + e_2 \otimes \theta^1 \wedge \theta^3).
    \end{aligned}
    \end{equation}

\medskip

    Now first suppose that the trace of the torsion tensor is non-vanishing. Then, taking its trace we find
    \begin{equation}
    \begin{aligned}
            V_\ai :=&  T^\bb_{~\bb \ai} = \left( \frac{ \omega_{1,u} m^2_{~2}-\omega_{2,u}m^2_{~1}}{m^1_{~1} m^2_{~2}} \right) \theta^2 + \left( \frac{\omega_{2,u}}{m^2_{~2}} \right) \theta^3 + (\ln m^1_{~1} m^2_{~2} )_{,u} \theta^1.
    \end{aligned}
    \end{equation}
    Comparing with $\mathcal{L}_{\ell} \omega$, we find that
    $$V_\ai - \mathcal{L}_{\ell} \omega_\ai = (\ln m^1_1 m^2_2)_{,u} \omega\,.$$
    Now suppose that $V(\ell) \neq 0$. Then, it follows that 
    \begin{equation}
        \tfrac{1}{V(\ell)} (V - \mathcal{L}_\ell \omega) = \omega\,.
    \end{equation}
    \noindent It thus follows that if
    $$\nabla_{(\ai} \gamma_{\bb)} =: N_{\ai \bb}\,,$$
    where
    $\gamma := \tfrac{1}{V(\ell)} (V - \mathcal{L}_\ell \omega)$, then $\nabla \odot \omega = N$, as required.

    On the other hand,
    if the trace of the torsion tensor is horizontal, meaning $V(\ell) = 0$, so that $V = \mathcal{L}_\ell \omega$, or the torsion tensor is trace-free, we may use the fact that the torsion tensor with all indices lowered can be written as 
    \begin{equation}
        T_{\ai \bb \cc} = (\mathcal{L}_\ell g_{\ai[\bb}) \omega_{\cc]}
    \end{equation}
    \noindent Taking the covariant derivative of this tensor and imposing the second condition in the theorem statement  yields $$\Gamma^1_{~22} = N_{22}, \quad \Gamma^1_{~33} = N_{33} \text{ and } \Gamma^1_{~32} = N_{23} +  \frac{\omega_{2,u} \omega_1- \omega_{1,u} \omega_2 - \omega_{2,x}+\omega_{1,y}}{2 m^1_{~1} m^2_{~2}}$$ from which it follows that $ \nabla \odot\omega = N$.

    In either case, it is easy to check by way of Equation~\ref{eqn:Thm32b} that $N$ is horizontal, meaning that $N(\ell,\cdot) = 0$.
    
    Proving the opposite direction is straightforward.

    \end{proof}

\end{document}
